\documentclass[11pt,a4paper]{article}

\usepackage[utf8]{inputenc}
\usepackage[T1]{fontenc}
\usepackage{lmodern}
\usepackage[french]{babel}
\usepackage[margin=2.1cm]{geometry}
\usepackage{microtype}
\usepackage{amsmath,amssymb}
\usepackage{eurosym}
\usepackage{booktabs}
\usepackage{array}
\usepackage[table]{xcolor}
\usepackage{graphicx}
\usepackage{caption}
\usepackage{enumitem}
\usepackage[most]{tcolorbox}

\definecolor{navy}{HTML}{1F3864}
\definecolor{lightblue}{HTML}{D6E0F0}
\definecolor{accent}{HTML}{C0491F}

\captionsetup{font=small,labelfont={bf,color=navy},labelsep=period}
\graphicspath{{../figures/}}

\newtcolorbox{cle}[1][]{colback=lightblue!40,colframe=navy,boxrule=0.6pt,
  arc=2pt,left=8pt,right=8pt,top=6pt,bottom=6pt,#1}

\emergencystretch=1.2em

\title{\vspace{-1.3cm}\color{navy}\bfseries L'usage des bases~: quelle source répond à quelle question\\[3pt]
  \large Fréquence et sévérité, ce que chacune peut porter et ce qu'elle ne peut pas}
\author{Kélian Kaddouri}
\date{3 août 2026}

\begin{document}
\maketitle
\vspace{-0.8cm}

\begin{cle}
\textbf{Le principe.} On ne cherche pas \emph{la} bonne base. Chaque source est utilisée
pour la seule question qu'elle est capable de trancher, et l'on déclare ce qui reste hors
d'atteinte. Cette note dit, brique par brique, quelle donnée alimente quel paramètre, à
quel niveau d'observation, et avec quelle marge d'erreur. Elle contient aussi la réponse à
la question posée sur la loi d'Erlang.
\end{cle}

% =====================================================================
\section*{1. Trois sources, trois rôles disjoints}

\begin{center}\small
\renewcommand{\arraystretch}{1.3}
\begin{tabular}{@{}p{3.3cm} p{4.6cm} p{7.4cm}@{}}
\toprule
\textbf{Source} & \textbf{Contenu} & \textbf{Ce qu'elle alimente} \\
\midrule
SAS OpRisk Global Data & pertes opérationnelles en dollars, catégories de Bâle,
$4\,329$ transitions annuelles &
\textbf{le niveau de sévérité} ($\xi=0{,}595$) \emph{et} \textbf{la fréquence d'entrée}.
C'est la source porteuse. \\
Privacy Rights Clearinghouse & $15\,053$ incidents, exprimés en \emph{nombre
d'enregistrements} &
\textbf{un axe de sensibilité} sur la sévérité, après conversion Jacobs
($\xi=1{,}033$, plafond $40$~M\euro). \\
7 rapports post-incident & $17$ enchaînements codés &
\textbf{le sens} des liens de contagion. Rien d'autre~: ni volume, ni amplitude. \\
\bottomrule
\end{tabular}
\end{center}

\medskip
Deux bases publiques ont par ailleurs été testées et \textbf{écartées}~: la Data Breach
Chronology (manque de volume, taxonomie par vecteur d'attaque et non par domaine de
contrôle, date piégée des deux côtés) et VCDB (le champ \emph{actor.Partner} passe de $56$
à $0$ sur la période, artefact de collecte).

% =====================================================================
\section*{2. La sévérité}

\subsection*{2.1. OpRisk, la source porteuse}

Périmètre~: secteur financier, sous-catégories \textsc{tic}, pertes converties en euros.
Le modèle n'ajuste que la queue, au dessus d'un seuil, par une loi de Pareto généralisée,
et conserve le corps empirique en dessous. L'ajustement retient environ $91$~excès au dessus
du seuil, pour $\xi=0{,}595$.

Ce chiffre est celui qui porte le capital principal, y compris le socle de $5\,275$~M\euro{}
et les bornes $[6\,858\,;\,8\,697]$~M\euro. \textbf{Il ne dépend d'aucune conversion}~: les
montants sont en dollars dès l'origine.

\subsection*{2.2. PRC, l'axe de sensibilité}

PRC ne contient pas de montants mais des \emph{nombres d'enregistrements compromis}. Le
passage aux euros se fait par la régression de Jacobs,
$\ln L = 7{,}68+0{,}76\ln X+\varepsilon$, dont le résidu a un écart-type de $0{,}523$.

Trois choses à savoir sur cette branche.

Elle donne $\xi=1{,}033$, donc au dessus de $1$~: la moyenne théorique de la sévérité
n'existe pas. Ce n'est pas une anomalie de calcul, c'est une propriété du périmètre.

La conversion était traitée comme \textbf{exacte}, en ignorant le résidu, ce qui écrasait la
dispersion. En le réinjectant, le niveau du $\Delta_{\text{DORA}}$ monte de $7\,\%$
($3\,600 \to 3\,861$~M\euro) et la queue s'épaissit ($\xi : 1{,}033 \to 1{,}072$), mais la
largeur relative de l'intervalle ne bouge pas ($1{,}43 \to 1{,}41$).

Et la conversion est \textbf{de sévérité seule}. Ce n'est pas un lien fréquence-sévérité~:
elle transforme un décompte en montant, la fréquence venant d'ailleurs.

% =====================================================================
\section*{3. La fréquence~: deux niveaux d'observation, deux paramètres}

C'est le point méthodologique de cette note, et celui où une erreur serait la plus coûteuse.

Le modèle de fréquence porte \textbf{deux} paramètres qui ne se lisent pas au même endroit~:
\begin{equation}
  N_j \mid Y \ \sim\ \mathcal{P}\bigl(\lambda_j\, m\bigr),
  \qquad m=\exp\!\Bigl(aY-\tfrac{a^2}{2}\Bigr),
  \qquad Y\sim\mathcal{N}(0,1).
  \label{eq:freq}
\end{equation}
Le taux $\lambda$ est une propriété d'\textbf{entité}. La charge systémique $a$ est un choc
\textbf{commun} à toutes les entités, tiré une fois par an.

\begin{cle}[colback=orange!8,colframe=orange!70!black]
\textbf{Pourquoi les mélanger serait une faute.} La surdispersion mesurée sur un panel
firme-année est dominée par l'\textbf{hétérogénéité de taille} des firmes, pas par le choc
commun~: le panel mélange des TPE et des banques systémiques. Le choc commun ne se lit
qu'au niveau \textbf{agrégé}, et seulement \textbf{après retrait de la tendance}, faute de
quoi on confond une dérive séculaire avec de la volatilité annuelle.
\end{cle}

\subsection*{3.1. Le taux par entité, lu sur le panel firme-année}

Fenêtre $2005$-$2022$. Au delà, les comptes s'effondrent ($34$ en $2023$, $39$ en $2024$
contre $70$ à $100$ auparavant)~: c'est la troncature à droite de la base, les événements
récents n'étant pas encore remontés, et non une baisse du risque.

\begin{center}\small
\begin{tabular}{@{}l r@{}}
\toprule
observations firme-année & $14\,944$ \;($5\,230$ firmes) \\
moyenne empirique & $0{,}0988$ événement \textsc{tic} / firme / an \\
variance empirique & $0{,}1179$ \\
indice de dispersion $\mathrm{Var}/\mathbb{E}$ & $1{,}19$ \\
part de zéros & $91{,}2\,\%$ \\
\midrule
Poisson & $\log L = -5\,023{,}8$ \\
Négative binomiale & $\log L = -4\,958{,}5$, \; forme $r=0{,}632$ \\
test du rapport de vraisemblance & $\mathrm{LR}=130{,}5$, \; $p=1{,}6\times10^{-30}$ \\
\bottomrule
\end{tabular}
\end{center}

\medskip
Le Poisson est rejeté sans ambiguïté. Et la sensibilité à la taille est forte~: sur les
seules firmes à dix événements ou plus, le taux passe de $0{,}099$ à $\mathbf{0{,}210}$, un
profil plus proche d'une entité effectivement soumise à DORA.

\subsection*{3.2. Le choc systémique, lu sur la série agrégée détrendée}

Sur la série sectorielle \textsc{tic} agrégée ($82$ événements par an en moyenne sur
$18$~ans), la tendance ajustée vaut $-2{,}4\,\%$ par an, et la dispersion \emph{résiduelle}
après retrait de tendance vaut $\varphi=2{,}24$, d'intervalle de confiance à $95\,\%$
$[1{,}24\,;\,5{,}18]$, large parce qu'il n'y a que $18$~points.

La charge implicite correspondante est $a=0{,}122$, d'intervalle $[0{,}054\,;\,0{,}223]$.
Sur le périmètre plus étroit cyber~$\times$~finance ($443$ événements, $24{,}6$ par an),
la même lecture donne $a=0{,}297$.

\begin{cle}
\textbf{Le verdict, et il est à assumer tel quel.} La valeur posée dans le modèle est
$A_{\text{LOAD}}=0{,}60$. Elle dépasse la charge empirique quel que soit le périmètre
($0{,}12$ ou $0{,}30$), et elle dépasse même la \textbf{borne haute} de l'intervalle de
confiance ($0{,}223$). Le calage n'est donc pas calibré, il est \textbf{délibérément
conservateur}, et cette conclusion résiste à l'incertitude d'estimation.
\end{cle}

% =====================================================================
\section*{4. La question de la loi d'Erlang}

La loi mélangeante ajustée sur le panel est une Gamma de forme $r=0{,}632$. Or~:

\begin{center}\small
\begin{tabular}{@{}l l@{}}
\toprule
Une \textbf{Erlang} & est une Gamma de forme \textbf{entière} $\ge 1$. \\
La forme ajustée & vaut $0{,}632$~: ni entière, ni supérieure à $1$. \\
\bottomrule
\end{tabular}
\end{center}

\medskip
L'Erlang la plus proche serait l'Erlang-1, c'est-à-dire un mélange exponentiel, mais ce
n'est qu'une \textbf{approximation}~: la Gamma réellement ajustée, de forme $0{,}63$, est
plus dissymétrique qu'une exponentielle.

\medskip
\noindent\textbf{À dire tel quel}~: le mélange Gamma est le bon objet, et l'Erlang n'en est
qu'un cas particulier à forme entière. Ici la forme n'est pas entière, donc parler d'Erlang
serait imprécis. La famille évoquée est la bonne, le nom ne l'est pas.

% =====================================================================
\section*{5. Les trois surdispersions, réconciliées}

Le modèle porte trois représentations de la surdispersion, longtemps coexistantes sans
avoir jamais été confrontées. Elles ne se contredisent pas, mais elles vivent à des
\textbf{échelles différentes}, et c'est cela qu'il faut savoir dire.

\begin{center}\small
\renewcommand{\arraystretch}{1.25}
\begin{tabular}{@{}l l r@{}}
\toprule
\textbf{Représentation} & \textbf{Échelle} & \textbf{Indice $\mathrm{Var}/\mathbb{E}$} \\
\midrule
mélange Gamma, $r=0{,}63$ & firme-année (micro, $\mu=0{,}099$) & $1{,}16$ \\
charge lognormale, $a=0{,}60$ & agrégée (portefeuille) & $10{,}34$ \\
facteur négatif binomial, $\varphi=9{,}20$ & agrégée (portefeuille) & $9{,}20$ \\
\bottomrule
\end{tabular}
\end{center}

\medskip
À l'échelle du moteur ($\lambda=21{,}6$ amorces par an, soit $582$ événements sur $27$ ans),
la correspondance entre les deux mesures agrégées se vérifie~: $\varphi=9{,}20$ implique
$a=0{,}568$ contre $0{,}60$ posé, et $a=0{,}60$ implique $\varphi=10{,}34$ contre $9{,}20$
posé. Elles \textbf{coïncident à $5\,\%$ près sur $a$} et à $12\,\%$ près sur $\varphi$~: ce
sont bien deux écritures de la même chose.

Le $r=0{,}63$ du panel, lui, vit à l'échelle d'une \emph{seule} firme et ne se compare pas
directement aux deux autres. C'est le piège de lecture principal de cette brique.

\medskip
\noindent\textbf{Et le $\varphi$ empirique est bien plus bas que le posé.} Sur la série
cyber~$\times$~finance détrendée, la dispersion résiduelle vaut $\varphi=3{,}27$, soit
$\varphi\approx3{,}0$ une fois remise à l'échelle du moteur, contre $9{,}20$ retenu. Le
calage est donc conservateur d'un facteur trois environ.

\subsection*{5.1. Ce que ce choix coûte, en capital}

\begin{center}\small
\begin{tabular}{@{}l r r@{}}
\toprule
\textbf{Loi de comptage} & \textbf{SCR (VaR $99{,}5\,\%$)} & \textbf{Moyenne} \\
\midrule
Poisson ($\varphi\to1$) & $7\,838$~M\euro & $1\,019$~M\euro \\
NB, $\varphi=3{,}0$ (empirique détrendé) & $7\,981$~M\euro & $1\,018$~M\euro \\
NB, $\varphi=9{,}20$ (posé) & $8\,553$~M\euro & $1\,054$~M\euro \\
NB, $\varphi=10{,}3$ (équivalent $a=0{,}60$) & $8\,534$~M\euro & $1\,028$~M\euro \\
\bottomrule
\end{tabular}
\end{center}

\medskip
Deux lectures. La \textbf{moyenne est invariante}, par construction, puisque
$\mathbb{E}[N]=\lambda$ dans les quatre lois~: ce que la surdispersion déplace, c'est la
VaR, pas l'espérance. Et l'effet reste \textbf{modéré}, $+9\,\%$ de Poisson au calage posé,
parce qu'à $99{,}5\,\%$ la queue reste largement portée par une perte dominante de sévérité.

Enfin, à variance \emph{égale}, la forme du mélange compte encore~: le quantile
$99{,}5\,\%$ du comptage vaut $74$ sous un mélange Gamma et $82$ sous un mélange lognormal,
soit $11\,\%$ d'écart. \textbf{Apparier la variance n'apparie pas la queue.}

% =====================================================================
\section*{6. Pourquoi la fréquence est calibrable et la direction ne l'est pas}

C'est la question qui revient, et la réponse est nette.

Calibrer la matrice de contagion $W$ exige trois choses simultanément~: un horodatage fin,
une taxonomie par domaine de contrôle, et la \textbf{direction}. Or la direction tombe sous
le test placebo ($z=-0{,}33$).

Calibrer une fréquence d'entrée n'exige \textbf{rien} de tout cela. Il suffit de savoir
qu'un événement \textsc{tic} est survenu chez telle entité telle année. Le pas annuel, fatal
pour $W$, suffit parfaitement pour un comptage.

\medskip
\noindent C'est un problème d'identification \textbf{strictement plus facile}, et c'est
pourquoi les deux briques n'ont pas le même statut~: la fréquence est confrontée à la
donnée, la direction est bornée.

% =====================================================================
\section*{7. Les biais à déclarer}

\textbf{Le seuil de collecte.} Le taux posé est de $12$ à $22$ incidents par an selon la
brique, contre $0{,}099$ à $0{,}210$ observés par firme-année. L'écart n'est pas une erreur
de calage~: la base ne retient que les pertes \textbf{au dessus d'un seuil de collecte}. On
calibre donc la fréquence des événements \textsc{tic} \emph{matériels}, générateurs de
perte, qui est la bonne maille pour un modèle de capital, et non la fréquence de tous les
incidents au sens du reporting DORA. Les deux nombres ne mesurent pas la même chose.

\textbf{Le biais de taille.} OpRisk sur-représente les grandes entités. Le taux de $0{,}210$
sur les firmes à dix événements ou plus est le plus représentatif d'une entité soumise à
DORA, et c'est celui qu'il faut citer.

\textbf{La troncature à droite.} Les deux dernières années de la base sont incomplètes. La
fenêtre s'arrête donc à $2022$, et il faut le dire avant qu'on ne s'étonne d'une baisse
apparente du risque.

\textbf{Le nombre de points.} La charge systémique repose sur $18$~observations annuelles.
Son intervalle de confiance est large, d'un facteur quatre. C'est aussi pourquoi le calage
conservateur est défendable~: il n'est pas contredit par la donnée, il la borne par le haut.

% =====================================================================
\begin{cle}
\textbf{En une phrase par brique.} La \textbf{sévérité} vient d'OpRisk, en dollars, sans
conversion, et PRC n'est qu'un axe de sensibilité. La \textbf{fréquence} se lit à deux
niveaux, le taux sur le panel firme-année et le choc commun sur la série agrégée détrendée,
et les mélanger serait une faute. Le calage retenu est \textbf{conservateur d'un facteur
trois} sur la surdispersion, ce qui est un choix de prudence et non une calibration, et il
déplace la VaR de $+9\,\%$ sans toucher la moyenne. La \textbf{direction}, elle, ne se
calibre sur aucune de ces bases, et c'est le seul paramètre du modèle qui reste borné plutôt
qu'estimé.
\end{cle}

\smallskip
\sloppy
\noindent\footnotesize\textcolor{navy}{\textbf{Sources.}} Fréquence d'entrée~:
\texttt{2\_donnees/08b\_calibration\_frequence\_entree.py}, figure
\texttt{N\_frequence\_entree.png}. Réconciliation des surdispersions et effet sur le SCR~:
\texttt{2\_donnees/41\_simulation\_nb\_comptage.py}, figure \texttt{N2\_simulation\_nb.png}.
Faisabilité et écartement des sources~:
\texttt{2\_donnees/05\_faisabilite\_donnees.py}. Résidu de la conversion Jacobs~:
\texttt{5\_etats/21\_prc\_jacobs\_residu.py}. Chapitre~5 du mémoire pour les données et
leurs limites.

\end{document}
